% Standalone problem document, generated from the adjacent README.md.
% Compile with XeLaTeX. Mathematical content is maintained in Markdown.
\documentclass[11pt,a4paper]{article}
\usepackage[fontsize=10.5pt]{scrextend}
\usepackage[margin=22mm,headheight=15pt,headsep=9mm,footskip=11mm]{geometry}
\usepackage{fontspec}
\setmainfont{texgyrepagella}[Extension=.otf,UprightFont=*-regular,BoldFont=*-bold,ItalicFont=*-italic,BoldItalicFont=*-bolditalic]
\setsansfont{texgyreheros}[Extension=.otf,UprightFont=*-regular,BoldFont=*-bold,ItalicFont=*-italic,BoldItalicFont=*-bolditalic]
\setmonofont{texgyrecursor}[Extension=.otf,UprightFont=*-regular,BoldFont=*-bold,ItalicFont=*-italic,BoldItalicFont=*-bolditalic,Scale=MatchLowercase]
\usepackage{amsmath,amssymb}
\usepackage{unicode-math}
\setmathfont{texgyrepagella-math.otf}
\usepackage{xcolor,microtype,enumitem,needspace,fancyhdr}
\usepackage{bookmark}
\definecolor{ink}{HTML}{17344B}
\definecolor{accent}{HTML}{197D80}
\definecolor{muted}{HTML}{596773}
\hypersetup{colorlinks=true,linkcolor=accent,urlcolor=accent,pdftitle={RA-20 - Critical-point
counts for symmetric rank-two approximation with diagonal
zeros},pdfauthor={Open Problems in Numerical Linear Algebra}}
\urlstyle{same}
\setlength{\parindent}{0pt}
\setlength{\parskip}{5pt plus 1pt minus 1pt}
\linespread{1.04}
\setlength{\emergencystretch}{3em}
\widowpenalty=10000
\clubpenalty=10000
\displaywidowpenalty=10000
\allowdisplaybreaks[1]
\setlist{nosep,leftmargin=1.5em}
\providecommand{\tightlist}{\setlength{\itemsep}{0pt}\setlength{\parskip}{0pt}}
\providecommand{\pandocbounded}[1]{#1}
\pagestyle{fancy}
\fancyhf{}
\fancyhead[L]{\sffamily\footnotesize\color{muted}OPEN PROBLEMS IN NUMERICAL LINEAR ALGEBRA}
\fancyhead[R]{\sffamily\bfseries\footnotesize\color{accent}RA-20}
\fancyfoot[L]{\parbox[b]{0.9\textwidth}{\sffamily\fontsize{8}{10}\selectfont\color{muted}Editorial ratings; a bounded search cannot certify that a problem remains unsolved.\\Literature check: 2026-09-11}}
\fancyfoot[R]{\sffamily\footnotesize\color{muted}\thepage}
\renewcommand{\headrulewidth}{0.3pt}
\renewcommand{\headrule}{\hbox to\headwidth{\color{accent}\leaders\hrule height \headrulewidth\hfill}}
\makeatletter
\renewcommand{\section}{\Needspace{5\baselineskip}\@startsection{section}{1}{0pt}{12pt plus 2pt minus 2pt}{4pt}{\sffamily\bfseries\large\color{ink}}}
\renewcommand{\subsection}{\Needspace{4\baselineskip}\@startsection{subsection}{2}{0pt}{9pt plus 1pt minus 1pt}{3pt}{\sffamily\bfseries\normalsize\color{ink}}}
\makeatother
\setcounter{secnumdepth}{0}
\begin{document}
{\sffamily\small\color{muted}Randomized and low-rank approximation\par}
\vspace{3pt}
{\raggedright\hyphenpenalty=10000\exhyphenpenalty=10000\sffamily\bfseries\fontsize{21}{24}\selectfont\color{ink}Critical-point
counts for symmetric rank-two approximation with diagonal zeros\par}
\vspace{7pt}
{\sffamily\small\textbf{Difficulty:} challenging\quad\textbf{Importance:} interesting
to specialist\par}
{\sffamily\small\bfseries\color{accent}Status: Solved\par}
\vspace{4pt}
{\color{accent}\hrule height 0.8pt}
\vspace{6pt}
\textbf{Topic:} Symmetric structured low-rank approximation

\textbf{Rating rationale:} Historical rating of the proposed all-order
formulas. The negative resolution below is elementary; it does not
establish corrected formulas in the remaining cases.

\subsection{Negative resolution -
2026-09-11}\label{negative-resolution---2026-09-11}

The displayed universal conjecture is false: \textbf{\(e_{3,3}=3\),
whereas its formula gives \(4\)}. For \(n=s=3\), write the off-diagonal
entries as \((a,b,c)\). The determinant is \(2abc\), so the variety is
the union of three coordinate planes. Each smooth component has exactly
one simple critical point for generic full-Frobenius data. Their
intersections are singular and are excluded by the definition below.

\href{https://github.com/ajt60gaibb/OpenProblemsInNLA/blob/main/randomized-and-low-rank-approximation/RA-20/solution.md}{Complete
proof} ·
\href{https://github.com/ajt60gaibb/OpenProblemsInNLA/blob/main/randomized-and-low-rank-approximation/RA-20/solution.pdf}{Proof
PDF} ·
\href{https://github.com/ajt60gaibb/OpenProblemsInNLA/blob/main/randomized-and-low-rank-approximation/RA-20/solution.tex}{Standalone
TeX} ·
\href{https://github.com/ajt60gaibb/OpenProblemsInNLA/blob/main/references/research-expansion-2026-09-11/ra20-resolution/README.md}{Independent
audit and exact verification}.

This counterexample was identified and independently checked during the
Codex maintainer audit. It is automated-agent verification, not external
human peer review or formal certification. The source's Table 7 repeats
the value \(4\); this is not a transcription error in the entry. The
complete original target, including its dimension range and all four
formulas, is retained below. No conclusion about corrected formulas or
the other parameter cases is claimed, and no priority claim is made.

\subsection{Statement}\label{statement}

For \(s\in\{1,2,3,4\}\) and \(n\ge\max\{3,s\}\), define

\[W_{n,s}=\{X\in\mathbb C^{n\times n}:X^{\mathsf T}=X,
\ \operatorname{rank}X\le2,\ x_{11}=\cdots=x_{ss}=0\}.\]

For generic symmetric \(U\in\mathbb C^{n\times n}\), let \(e_{n,s}\) be
the number of complex critical points on the smooth locus of \(W_{n,s}\)
of

\[d_U(X)=\sum_i(x_{ii}-u_{ii})^2+
2\sum_{i<j}(x_{ij}-u_{ij})^2.\]

Generic means outside a proper algebraic exceptional set. A point is
critical if the differential vanishes on its tangent space. Thus
\(e_{n,s}\) is the Euclidean distance degree for the bilinear extension
of the \textbf{full Frobenius metric}; the off-diagonal terms have
weight two, and complex conjugation is absent.

\textbf{Conjecture (Kubjas--Sodomaco--Tsigaridas, Conjecture 5.6,
\(n\ge3\)).}

\[e_{n,s}=\begin{cases}
3(n-1)-2,&s=1,\\
9(n-2)-2,&s=2,\\
27(n-3)+4,&s=3,\\
81(n-4)+28,&s=4.
\end{cases}\]

All four zero counts form one target. The explicit \(n\ge3\) restriction
avoids a degenerate endpoint in the printed statement: when \(n=2\), the
rank bound is vacuous and both possible varieties are linear with ED
degree one, whereas the displayed \(s=2\) formula does not apply.

\subsection{Evidence and numerical
significance}\label{evidence-and-numerical-significance}

The source's Table 7 reports values matching its formulas through order
ten, but the \(n=s=3\) value conflicts with the exact calculation above.
The table therefore cannot establish the conjecture's validity. The
problem counts stationary candidates for symmetric Frobenius
approximation with prescribed diagonal zeros. It concerns fixed rank
two, unlike
\href{https://github.com/ajt60gaibb/OpenProblemsInNLA/blob/main/randomized-and-low-rank-approximation/RA-19/README.md}{corank-one
approximation in general square matrices}.

\subsection{References and status
check}\label{references-and-status-check}

\begin{itemize}
\tightlist
\item
  K. Kubjas, L. Sodomaco and E. Tsigaridas, \emph{Exact solutions in
  low-rank approximation with zeros}, Linear Algebra and its
  Applications 641 (2022), 67--97.
  \href{https://doi.org/10.1016/j.laa.2022.01.021}{DOI};
  \href{https://arxiv.org/abs/2010.15636v2}{current author manuscript},
  29 January 2022. Conjecture 5.6 and Table 7, manuscript p.21; §2 and
  §5 supply the distance convention.
\end{itemize}

The initial literature search on 2026-09-11 found no later resolution
and compared the formulas with Table 7. The subsequent independent
maintainer audit supplied the counterexample above, superseding the
initial Open classification. Excluding \(n=2\) does not remove the
admissible counterexample \(n=s=3\). The separate nonsymmetric formulas
in Conjecture 5.2 have a table/label discrepancy and are not imported
here. The permanent ID, original target and canonical path are retained.
\end{document}
